\chapter{Serial Interfaces: UART, SPI, I$^2$C}
\label{ch:uart}

A chip has a limited number of pins, but there is a lot of data to transfer. If
you send a byte over eight wires at once (a \term{parallel} interface), you
need eight pins and eight traces on the board, and at high frequencies the
wires start to interfere with each other. That is why data between chips is
most often transferred \term{serially}: over a single wire, bit by bit. In
this chapter we will look at the three most common serial interfaces and build
one of them---UART---completely on an FPGA, from the transmitter to the
receiver.

\section{Three Interfaces}

\begin{table}[H]
\centering
\caption{Comparison of serial interfaces}
\small
\begin{tabularx}{\textwidth}{l X X X}
\toprule
 & \textbf{UART} & \textbf{SPI} & \textbf{I$^2$C} \\
\midrule
Wires & 2 (TX, RX) + ground & 4 (SCK, MOSI, MISO, CS) & 2 (SDA, SCL) \\
Clock & none, the speed is agreed in advance & yes, the master generates SCK & yes, the master generates SCL \\
Devices & 2 (point-to-point) & 1 master and several slaves (selected by CS lines) & up to 127 addresses on 2 wires \\
Speed & up to a few Mbit/s & tens of Mbit/s & 0.1--3.4 Mbit/s \\
Typical uses & console, GPS, Bluetooth modules & Flash, displays, ADCs & sensors, clocks, EEPROM \\
\bottomrule
\end{tabularx}
\end{table}

\begin{diybox}[How many wires per interface]
\textbf{Task.} Compare the three interfaces using ready-made examples and make sure the number of signals matches the table above.

\textbf{In Logra.} In the From an example card, open three projects: UART transmitter 8N1, SPI mode 0: internal loopback, and I²C: START and STOP. On the Code tab, look at the port list of the top-level module; on the Pins tab, see which of them the platform needs.

\textbf{Expected result.} UART has separate transmit and receive lines (\code{tx}, \code{rx}) and no clock; SPI has \code{SCK}, \code{MOSI}, \code{MISO}, \code{CS}; I$^2$C has a single pair of lines, \code{SDA} and \code{SCL}.

\textbf{How to check.} Make a table ``interface---signals---has a clock'' and compare it with the table above; in the SPI example, note that the exchange is internal (loopback) and the board pins are not used.

\stars{1} repeat: fill in the table for the three projects. \stars{2} modify: choose an interface for a sensor that has only one free pin left, and justify your choice. \stars{3} project: Security Panel stage---choosing the communication channel; acceptance criterion: one paragraph explains why the panel uses UART to talk to the computer.
\end{diybox}

\section{UART}

\term{UART} (universal asynchronous receiver-transmitter) is the simplest
interface; it is present in every microcontroller and in the USB bridge on your
board (it is exactly what connects the FPGA to the computer over USB).
Asynchronous means that there is \emph{no shared clock signal}: the transmitter
and the receiver agree on the speed (the baud rate) in advance, and the
receiver adjusts itself to the start of each byte.

At rest, the line is held at 1. A byte is sent as a frame: a start bit (0),
eight data bits \emph{least significant bit first}, and a stop bit (1). This
format is called 8N1 (8 data bits, no parity, 1 stop bit).

\begin{figure}[H]
\centering
\begin{tikzpicture}[x=1.0cm, y=0.8cm, font=\sffamily\scriptsize]
  \draw[FpgaBlue, very thick]
    (0,1) -- (1,1) -- (1,0) -- (2,0) -- (2,1) -- (3,1) -- (3,0) -- (8,0) -- (8,1) -- (9,1) -- (9,0) -- (10,0) -- (10,1) -- (11.5,1);
  \foreach \x/\t in {1.5/start, 2.5/{$d_0$}, 3.5/{$d_1$}, 4.5/{$d_2$}, 5.5/{$d_3$}, 6.5/{$d_4$}, 7.5/{$d_5$}, 8.5/{$d_6$}, 9.5/{$d_7$}, 10.5/stop}
    \node at (\x,1.7) {\t};
  \foreach \x/\t in {2.5/1, 3.5/0, 4.5/0, 5.5/0, 6.5/0, 7.5/0, 8.5/1, 9.5/0}
    \node[text=FpgaGray] at (\x,-0.5) {\t};
  \foreach \x in {1,2,...,11} \draw[gray!50, dashed] (\x,-0.2) -- (\x,1.4);
  \node[anchor=east] at (0,1) {1};
  \node[anchor=east] at (0,0) {0};
  \node[text=FpgaGray] at (0.5,0.5) {idle};
  \draw[<->] (1,-1.1) -- node[below] {$T_{\text{bit}}$} (2,-1.1);
  \node[font=\sffamily\small] at (6,-2.0) {Frame for the letter A = \code{0x41} = \code{01000001$_2$}: the least significant bit goes first};
\end{tikzpicture}
\caption{A UART 8N1 frame}
\label{fig:uart_frame}
\end{figure}

\subsection{Baud Rate}

The duration of one bit, $T_{\text{bit}}$, equals $1/\text{baud rate}$. For
115\,200 baud it is about $8.7$~$\mu$s. The FPGA runs at 27~MHz, so one bit
takes
\[
\text{CLKS\_PER\_BIT} = \frac{27\,000\,000}{115\,200} \approx 234 \text{ clock cycles}.
\]
Both the transmitter and the receiver count clock cycles and change (or read)
a bit every \code{CLKS\_PER\_BIT} cycles.

\subsection{The Transmitter}

The transmitter is simpler: it stores the assembled frame in a shift register
and puts its bits on the line on a timer. Between bytes, the line stays at
rest (1). The state machine has two states: \code{IDLE}, where it waits for the
\code{start} signal, and \code{SEND}, where the ten bits of the frame are sent
(start, 8 data bits, stop).

\subsection{The Receiver}

The receiver is more complicated because it does not know when a frame will
begin, and it has no shared clock. It works like this:

\begin{enumerate}
  \item It waits until the line drops to 0: this is the start of the start bit.
  \item It counts \emph{half a bit} to land in the middle of the start bit and
        checks that the line is still at zero (otherwise it was noise).
  \item After that, it reads the line exactly once per bit---always in the middle
        of the bit, where the signal has settled and is farthest from the edges.
  \item It reads the stop bit: it must be a 1, otherwise the frame is considered
        corrupted.
\end{enumerate}

In addition, the \code{rx} input comes from outside and has no relation to the
FPGA clock, so there are two flip-flops in a row at the receiver input (a
synchronizer, see Chapter~\ref{ch:clocks}).

\vfile{\codepath/08_uart/uart.v}{code/08\_uart/uart.v}

The file contains three modules. \code{uart\_tx} and \code{uart\_rx} are the
transceivers themselves, with the \code{CLKS\_PER\_BIT} parameter, and
\code{uart\_demo} is their demonstration: the transmitter keeps sending the
letters from A to Z, and the \code{tx} output is connected to the receiver
input of the same module, so the receiver gets what was sent. For viewing in
the simulator, the speed is made very high (\code{CLKS\_PER\_BIT = 8});
otherwise you would have to wait for thousands of clock cycles.

\shot{n-uart-wave}{Sending the letter A in Logra: first the start bit and eight data bits appear on \code{tx}, then the receiver raises \code{rx\_valid} and \code{0x41} appears in \code{rx\_data}}{fig:n-uart-wave}

On the waveform (Fig.~\ref{fig:n-uart-wave}) you can count the frame: the line
goes down (start), several bits are sent, and at the end, exactly 10 bit
periods later, the receiver announces a ready byte with a \code{rx\_valid}
pulse and sets \code{rx\_data}.

\begin{infobox}[How to connect to a computer]
The USB bridge on your board brings UART out to the computer. All you need is
to set \code{CLKS\_PER\_BIT = 234} (115\,200 baud at 27~MHz), assign the
\code{tx} and \code{rx} pins in the pin planner, load the design, and open the
UART terminal tab in the Connect section: the letters you send will appear
there.
\end{infobox}

\begin{diybox}[Send the letter A]
\textbf{Task.} Send the letter A (\code{0x41}) over UART, break down the frame using the waveform, and pass it through the receiver.

\textbf{In Logra.} Open UART transmitter 8N1 and UART receiver 8N1, or \code{uart.v} from the book (\code{uart\_demo}, \code{CLKS\_PER\_BIT = 8}). Waveforms tab: the signals \code{tx}, \code{rx\_valid}, \code{rx\_data}.

\textbf{Expected result.} The line goes down for the start bit, then come the bits $1,0,0,0,0,0,1,0$ (least significant first), then the stop 1; the receiver raises \code{rx\_valid} and sets \code{rx\_data = 0x41}.

\textbf{How to check.} Waveforms (Fig.~\ref{fig:n-uart-wave}): count the bits of the frame, make sure the width of a bit equals \code{CLKS\_PER\_BIT} clock cycles and that \code{rx\_valid} arrives at the end of the frame. For the board: set \code{CLKS\_PER\_BIT = 234}, assign the pins on the Pins tab, then Connect tab $\to$ UART terminal.

\stars{1} repeat: the letter A on the waveform. \stars{2} modify: send \code{0x55} and make sure the line alternates strictly. \stars{3} project: Security Panel stage---sending the state over UART; acceptance criterion: at every change of \code{state}, a byte (0, 1, or 2) appears in the terminal at 115\,200 baud.
\end{diybox}

\section{SPI}

\term{SPI} has a shared clock, \code{SCK}, which makes it simpler and faster.
The master generates the clock signal and the data on the \code{MOSI} line
(master out, slave in), and at the same time receives data from the slave on
the \code{MISO} line. The slave is selected by the active-low \code{CS} (chip
select) signal.

\begin{figure}[H]
\centering
\begin{tikzpicture}[x=0.7cm, y=0.8cm, font=\sffamily\scriptsize]
  \node[anchor=east] at (-0.2,2.7) {CS};
  \draw[FpgaBlue, very thick] (0,3) -- (1,3) -- (1,2.4) -- (17,2.4) -- (17,3) -- (18,3);
  \node[anchor=east] at (-0.2,1.5) {SCK};
  \draw[FpgaBlue, very thick] (0,1.2) -- (2,1.2);
  \foreach \i in {0,...,7}
    \draw[FpgaBlue, very thick] (2+\i*2,1.2) -- (2+\i*2,1.8) -- (3+\i*2,1.8) -- (3+\i*2,1.2) -- (4+\i*2,1.2);
  \node[anchor=east] at (-0.2,0) {MOSI};
  \draw[FpgaGreen!60!black, very thick] (0,-0.3) -- (1,-0.3) -- (1,0.3) -- (3,0.3) -- (3,-0.3) -- (5,-0.3) -- (5,0.3) -- (9,0.3) -- (9,-0.3) -- (13,-0.3) -- (13,0.3) -- (15,0.3) -- (15,-0.3) -- (17,-0.3);
  \foreach \i/\v in {0/1, 1/0, 2/1, 3/1, 4/0, 5/0, 6/1, 7/0} \node at (2+\i*2,-1.0) {\v};
  \foreach \i in {0,...,7} \draw[gray!50, dashed] (2+\i*2,-0.5) -- (2+\i*2,1.9);
  \node[align=left, anchor=west] at (0,-2.1) {Data changes on the falling edge of SCK and is read by the slave on the rising edge (dashed), mode 0};
\end{tikzpicture}
\caption{Sending the byte \code{10110010} over SPI. Data changes on the falling clock edge and is read on the rising edge}
\end{figure}

SPI has four \term{modes}---combinations of the clock polarity at rest (CPOL)
and the phase (CPHA, reading on the first or the second edge). The modes must
be checked against the slave device's documentation. The simplest SPI
transmitter in an FPGA is the shift register from Chapter~\ref{ch:regs}, which
shifts on every SCK cycle, plus a frequency divider that generates SCK itself.

\begin{diybox}[SPI: shifting and sampling]
\textbf{Task.} Send the byte \code{0xA6} over SPI in mode 0 and check that the master received the same thing through the internal loopback.

\textbf{In Logra.} The SPI mode 0: internal loopback project (From an example card), parameter \code{HALF\_TICKS = 1}. Waveforms tab: \code{sck}, \code{mosi}, \code{cs\_n}, \code{received}.

\textbf{Expected result.} \code{SCK} is 0 at rest; the data on \code{MOSI} changes on the falling edge and is read on the rising edge; after eight clock cycles, \code{received = 0xA6}.

\textbf{How to check.} Waveforms: compare the order of the \code{MOSI} bits with \code{0xA6 = 10100110$_2$} (most significant first) and with the SPI frame figure.

\stars{1} repeat: \code{0xA6}. \stars{2} modify: send \code{0xC3} and predict the value of \code{received}. \stars{3} project: Security Panel stage---status byte over SPI; acceptance criterion: \code{cs\_n} is low for exactly 8 cycles of \code{SCK}, and the data does not change on the rising edge.
\end{diybox}

\section{I$^2$C}

\term{I$^2$C} uses only two wires, but in return it makes the protocol more
complex. \code{SDA} carries data and \code{SCL} carries the clock. Both lines
are open-drain and pulled up to the supply by resistors
(Chapter~\ref{ch:electrical}), so any device can pull the line down to 0, while
the 1 is created by the pull-up.

A transfer begins with the \term{START} condition (SDA falls while SCL is
high). Then come the 7-bit device address, the direction bit (read or write),
and the \term{ACK} acknowledgment: the addressed chip itself pulls SDA low on
the ninth clock cycle. Then bytes are transferred, each with an
acknowledgment, and at the end comes the \term{STOP} condition (SDA rises while
SCL is high). Many devices can sit on one line, each with its own address, and
collisions do not happen precisely thanks to the open drain.

An I$^2$C project is more complex: you need a state machine with dozens of
states, and a real device acknowledges reception. The Logra lab has models of
an OLED display, a clock, and a memory on I$^2$C that respond bit by bit, so you
can test your state machine before turning on the board.

\begin{diybox}[I\texorpdfstring{$^2$}{2}C: START and STOP]
\textbf{Task.} Find the START and STOP conditions on the waveform and make sure an open-drain line cannot be forcibly pulled up to 1.

\textbf{In Logra.} The I²C: START and STOP project (From an example card), \code{STEP\_TICKS = 1}; Waveforms tab: \code{sda}, \code{scl}, \code{drive\_sda\_low}.

\textbf{Expected result.} START: \code{SDA} falls while \code{SCL} is high; then one training bit; STOP: \code{SDA} rises while \code{SCL} is high. There is no address or ACK in this example.

\textbf{How to check.} Waveforms: mark both conditions. Then open the Lab tab (the Check group), find the I$^2$C device in the parts list, and see how it responds to an address.

\stars{1} repeat: find START and STOP. \stars{2} modify: add the 7 address bits and an ACK bit to the sequence. \stars{3} project: Security Panel stage---an indicator on I$^2$C; acceptance criterion: the waveform shows START, the address, the direction bit, ACK, the data byte, and STOP.
\end{diybox}

\section{Questions and Problems}
\label{sec:tasks-uart}

The answers are collected in Appendix~\ref{app:answers}.

\begin{taskbox}[Chapter problems]
\begin{enumerate}
  \item \stars{1}~What \code{CLKS\_PER\_BIT} is needed for 9600 baud at a frequency of 27~MHz?
  \item \stars{1}~How long does it take to send one byte in 8N1 format at
        115\,200 baud? How many bytes per second can be sent?
  \item \stars{1}~Draw the UART frame for the byte \code{0x55} (\code{01010101}).
  \item \stars{2}~Why does a UART receiver read bits in the \emph{middle} of their duration rather than at the start?
  \item \stars{2}~The receiver's clock frequency is 3\,\% higher than the transmitter's. Will the
        reception of a 10-bit frame fail? What speed mismatch is acceptable?
  \item \stars{2}~Why is SPI faster than UART, and what is its drawback?
  \item \stars{2}~Find the bug. A UART receiver, after the line drops, starts reading bits immediately, counting
        \code{CLKS\_PER\_BIT} cycles from the first edge. At 115\,200 baud with a good
        cable everything works, but on a long wire failures appear. Why? How can it be fixed?
  \item \stars{3}~Project. A module sends the panel's status byte (0, 1, or 2) every second
        at 115\,200 baud. Acceptance criterion: on the waveform, the frame consists of
        a start bit, 8 bits least significant first, and a stop bit, and in the terminal a character arrives once per second.
  \item \stars{3}~Project. Connect the two modules \code{uart\_tx} and \code{uart\_rx} in one project
        and show the received byte on the LEDs. Acceptance criterion: all 256 values pass
        without corruption, and applying a wrong stop bit does not produce \code{rx\_valid}.
\end{enumerate}
\end{taskbox}
